\documentclass[aspectratio=169]{beamer}
\usetheme{Madrid}
\usecolortheme{default}
% --- Organization slides (TOC + section dividers) ---
\AtBeginSection[]
{
  \begin{frame}{Outline}
    \tableofcontents[currentsection]
  \end{frame}
}

\title{Lecture 3: Improving Performance}
\subtitle{From a Working Prototype to Reliable Results}
\author{University of Chicago}
\date{\today}

\begin{document}

\frame{\titlepage}

\section{Where We Are}

\begin{frame}{Review: What We've Covered}
\textbf{Lecture 1: Foundations}
\begin{itemize}
    \item Basic AI terminology (LLMs, tokens, context windows, APIs)
    \item How to work with these models programmatically
    \item Understanding pricing and token economics
\end{itemize}
\vspace{0.3cm}

\textbf{Lecture 2: Building AI Systems}
\begin{itemize}
    \item Vertical slices -- prove end-to-end flow works first
    \item Crawl, Walk, Run -- start simple, add complexity incrementally
    \item Built a resume scorer: resume + job req $\rightarrow$ 0--100 score
\end{itemize}
\end{frame}

\begin{frame}{Where We're Going}
\textbf{Last time}: We built a working resume scorer
\begin{itemize}
    \item Single prompt, single call, returns a score
    \item It works! But how \textit{well} does it work?
\end{itemize}
\vspace{0.3cm}

\textbf{Today}: Techniques to improve AI system performance
\begin{itemize}
    \item Walk through common failure modes
    \item Learn a toolkit of solutions
    \item Apply them to the resume scorer in the notebook
\end{itemize}
\vspace{0.3cm}

\textbf{Framework}: Problem $\rightarrow$ Why it happens $\rightarrow$ Solution
\end{frame}

\section{The Single-Shot Baseline}

\begin{frame}{Our Starting Point}
\textbf{The resume scorer from Lecture 2}:
\vspace{0.3cm}

\begin{center}
\fbox{Resume Text} + \fbox{Job Requirements}
$\;\longrightarrow\;$
\fbox{LLM: ``Score 0--100''}
$\;\longrightarrow\;$
\fbox{Score: 72}
\end{center}
\vspace{0.3cm}

\textbf{This is a vertical slice} -- it works end-to-end.\\
Now: how do we make it \textit{good}?
\end{frame}

\begin{frame}{What Does ``Good'' Look Like?}
\vspace{0.3cm}

\begin{itemize}
\item What are some characteristics of good? 
\item (Open Question)
\end{itemize}
\vspace{0.3cm}

\end{frame}

\begin{frame}{What Does ``Good'' Look Like?}
\vspace{0.3cm}

\begin{itemize}
\item \textbf{Accurate} -- scores reflect actual candidate quality
\item \textbf{Consistent} -- same resume gets the same score every time
\item \textbf{Differentiated} -- strong candidates score higher than weak ones
\item \textbf{Explainable} -- you can justify why a candidate got their score
\end{itemize}
\vspace{0.3cm}

\end{frame}

\begin{frame}{Model Iteration}
\textbf{Best practice}: Create a measurable result, then iterate
\vspace{0.3cm}

\begin{enumerate}
\item \textbf{Label some data} -- identify resumes where you know the answer
\item \textbf{Run your prompt} -- see how the model performs
\item \textbf{Identify failures} -- where does it get things wrong?
\item \textbf{Adjust and re-run} -- change the prompt, check if results improve
\end{enumerate}
\vspace{0.3cm}

\textbf{This is the loop we'll follow today}: identify a problem, apply a fix, measure improvement.
\end{frame}

\section{Common Problems and Solutions}

\section{Problem 1: Score Averaging}

\begin{frame}{Problem: Everything Gets a Similar Score}

\textbf{You run the scorer on 6 resumes and get}:
\vspace{0.3cm}

\begin{center}
\begin{tabular}{l c}
\textbf{Resume} & \textbf{Score} \\
\hline
Resume A & 74 \\
Resume B & 71 \\
Resume C & 68 \\
Resume D & 72 \\
Resume E & 66 \\
Resume F & 70 \\
\end{tabular}
\end{center}
\vspace{0.3cm}

\textbf{Everything is between 66--74.} The scorer can't tell the difference between a great candidate and a poor one.
\end{frame}

\begin{frame}{Why Score Bunching Happens}
\textbf{LLMs default to ``safe'' middle-range answers}:
\begin{itemize}
    \item Training data is full of moderate, hedging language
    \item Without strong guidance, models avoid extremes
    \item A vague prompt like ``score 0--100'' gives the model no anchor points
\end{itemize}
\vspace{0.3cm}

\textbf{The model doesn't know what you mean by ``good''}:
\begin{itemize}
    \item Is 5 years of Python experience good or average?
    \item Is a Master's degree required or a nice-to-have?
    \item Should missing a skill be a small or large penalty?
\end{itemize}
\vspace{0.3cm}

\textbf{Without criteria, the model invents its own} -- and they're inconsistent.
\end{frame}

\begin{frame}{Solution: Domain Context and Grounding}
\textbf{Grounding} = anchoring the model's output in specific facts, criteria, or source material rather than letting it generate from general patterns.
\vspace{0.3cm}

\textbf{Tell the model exactly what matters and how much}:
\vspace{0.2cm}

Instead of: ``Score this resume 0--100''
\vspace{0.2cm}

Provide:
\begin{itemize}
    \item \textbf{Specific criteria} with weights or importance levels
    \item \textbf{What ``high'' and ``low'' mean} for this role
    \item \textbf{Must-haves vs. nice-to-haves}
    \item \textbf{Anchor points}: ``A score of 90+ means the candidate meets all requirements. Below 40 means missing critical skills.''
\end{itemize}
\vspace{0.3cm}

\begin{center}
\fcolorbox{blue!80}{blue!10}{%
\begin{minipage}{0.85\textwidth}
\centering
\textbf{The more specific your criteria, the more differentiated the scores}
\end{minipage}
}
\end{center}
\end{frame}

\section{Problem 2: Hallucinated Reasoning}

\begin{frame}{Problem: The Model Makes Things Up}
\textbf{Hallucination} = when the model generates confident, plausible-sounding content that is not supported by the input data.
\vspace{0.3cm}

\textbf{You ask the model to explain its score}:
\vspace{0.2cm}

\begin{quote}
``The candidate has 8 years of Python experience and demonstrated strong leadership at a Fortune 500 company...''
\end{quote}
\vspace{0.2cm}

\textbf{But the resume says}: 3 years of Python, worked at a 20-person startup.
\vspace{0.3cm}

\textbf{The model generated a plausible-sounding justification that doesn't match the actual resume.}
\end{frame}

\begin{frame}{Why Hallucination Happens}
\textbf{LLMs are pattern completers, not fact checkers}:
\begin{itemize}
    \item They generate text that \textit{sounds right} based on patterns
    \item ``Senior developer'' in training data often correlates with ``8+ years''
    \item Without explicit instruction to reference the source, they fill gaps
\end{itemize}
\vspace{0.3cm}

\textbf{This is especially dangerous for scoring}:
\begin{itemize}
    \item Score might be based on hallucinated qualifications
    \item You can't trust the reasoning to explain the score
    \item Impossible to audit or defend decisions
\end{itemize}
\end{frame}

\begin{frame}{Solution: Require Evidence (Grounding)}
\textbf{Force the model to cite its sources}:
\vspace{0.3cm}

\begin{itemize}
    \item Require \textbf{exact quotes} from the resume for every claim
    \item Ask for evidence \textit{before} the conclusion
    \item Include an ``evidence'' field in structured output
    \item Validate that cited text actually appears in the resume
\end{itemize}
\vspace{0.3cm}

\textbf{Result}: Instead of inventing ``8 years of Python,'' the model must point to the actual text: ``Python developer, 2021--2024''
\vspace{0.3cm}

\begin{center}
\fcolorbox{green!80}{green!10}{%
\begin{minipage}{0.85\textwidth}
\centering
\textbf{Grounding turns the LLM from a ``creative writer'' into an ``evidence-based analyst''}
\end{minipage}
}
\end{center}
\end{frame}

\section{Problem 3: Inconsistency}

\begin{frame}{Problem: Same Resume, Different Scores}
\textbf{Run the same resume through the scorer three times}:
\vspace{0.3cm}

\begin{center}
\begin{tabular}{l c}
\textbf{Run} & \textbf{Score} \\
\hline
Run 1 & 78 \\
Run 2 & 65 \\
Run 3 & 82 \\
\end{tabular}
\end{center}
\vspace{0.3cm}

\textbf{A 17-point spread on the same input.}
\vspace{0.2cm}

\begin{itemize}
    \item Is this candidate a 65 or an 82?
    \item Would you trust this system to rank candidates?
    \item What if the hiring decision depends on a threshold?
\end{itemize}
\end{frame}

\begin{frame}{Why Inconsistency Happens}
\textbf{LLMs are non-deterministic by design}:
\begin{itemize}
    \item Temperature controls randomness in token selection
    \item Even at low temperature, outputs can vary
    \item Vague instructions amplify variance -- many ``correct'' answers
\end{itemize}
\vspace{0.3cm}

\textbf{The less constrained the task, the more variance you get}:
\begin{itemize}
    \item ``Score 0--100'' has enormous output space
    \item ``Classify as: Strong Match / Partial Match / No Match'' has 3 options
    \item More structure = less room for random drift
\end{itemize}
\end{frame}

\begin{frame}{Solution: Temperature and Structured Output}
\textbf{Reduce randomness and constrain the output}:
\vspace{0.3cm}

\begin{itemize}
    \item \textbf{Lower temperature} (0.0--0.3) for classification and scoring tasks
    \begin{itemize}
        \item Less creative, more consistent
        \item Save high temperature for brainstorming, not evaluation
    \end{itemize}
    \vspace{0.2cm}
    \item \textbf{Structured output schemas} force specific fields and types
    \begin{itemize}
        \item JSON schema with required fields
        \item Constrained ranges (score must be 0--100)
        \item Enumerated categories instead of free text
    \end{itemize}
    \vspace{0.2cm}
    \item \textbf{Specific criteria} reduce subjective interpretation
    \begin{itemize}
        \item ``Does the candidate have 5+ years Python?'' has less variance than ``Rate their Python skills''
    \end{itemize}
\end{itemize}
\end{frame}

\section{Problem 4: Missing Nuance}

\begin{frame}{Problem: The Scorer Ignores Important Factors}
\textbf{A strong candidate gets a mediocre score because}:
\begin{itemize}
    \item The scorer focused on keyword matching
    \item Missed that the candidate led a relevant project
    \item Didn't weigh industry experience appropriately
    \item Treated a PhD the same as a Bachelor's
\end{itemize}
\vspace{0.3cm}

\textbf{A weaker candidate scores nearly as high because}:
\begin{itemize}
    \item Had the right keywords sprinkled throughout
    \item Model couldn't distinguish depth from breadth
    \item Single prompt tries to evaluate everything at once
\end{itemize}
\vspace{0.3cm}

\textbf{One prompt asking ``score this resume'' is too much at once.}
\end{frame}

\begin{frame}{Solution: Decomposition}
\textbf{Break the scoring into focused sub-tasks}:
\vspace{0.2cm}

\begin{center}
\begin{tabular}{ccc}
 & \fbox{Extract Experience} & \\
\fbox{Resume} $\;\longrightarrow\;$ & \fbox{Extract Skills} & $\;\longrightarrow\;$ \fbox{Combine \& Score (in code)} \\
 & \fbox{Extract Education} & \\
\end{tabular}
\end{center}
\vspace{0.2cm}

\textbf{Each sub-task is simpler and more focused}:
\begin{itemize}
    \item Easier for the model to get right
    \item Easier for you to debug (which step went wrong?)
    \item Can use different prompts/models for different steps
\end{itemize}
\end{frame}

\begin{frame}{Decomposition: Key Idea}
\textbf{Use LLMs for what they're good at}:
\begin{itemize}
    \item Extracting information from unstructured text
    \item Classification and categorization
    \item Summarization
\end{itemize}
\vspace{0.3cm}

\textbf{Use code for what code is good at}:
\begin{itemize}
    \item Conditional logic (if/then rules)
    \item Math and scoring formulas
    \item Combining results from multiple steps
    \item Enforcing business rules consistently
\end{itemize}
\vspace{0.3cm}

\begin{center}
\fcolorbox{blue!80}{blue!10}{%
\begin{minipage}{0.85\textwidth}
\centering
\textbf{LLMs extract and classify. Code decides and calculates.}
\end{minipage}
}
\end{center}
\end{frame}

\section{Problem 5: Edge Cases}

\begin{frame}{Problem: Ambiguous Resumes Get Random Scores}
\textbf{Some resumes are genuinely hard to evaluate}:
\begin{itemize}
    \item Career changer with transferable but non-obvious skills
    \item Experience in adjacent technology (Java listed, role needs Python)
    \item Years of experience in a broader category but not the specific area
\end{itemize}
\vspace{0.3cm}

\textbf{Without guidance, the model handles these inconsistently}:
\begin{itemize}
    \item Sometimes gives benefit of the doubt, sometimes doesn't
    \item No consistent policy for ``close but not exact'' matches
    \item Each run may interpret the ambiguity differently
\end{itemize}
\end{frame}

\begin{frame}{Solution: Few-Shot Examples}
\textbf{Show the model how you want edge cases handled}:
\vspace{0.3cm}

\textbf{Include 2--5 examples in your prompt that cover}:
\begin{enumerate}
    \item \textbf{A clear strong match} -- what does ``high score'' look like?
    \item \textbf{A clear weak match} -- what earns a low score?
    \item \textbf{An ambiguous case} -- how should the model handle it?
\end{enumerate}
\vspace{0.3cm}

\textbf{Examples are more powerful than instructions}:
\begin{itemize}
    \item ``Score based on relevant experience'' is vague
    \item Showing a scored example of ``5 years Java $\rightarrow$ 60 for a Python role'' is concrete
    \item The model learns your \textit{policy}, not just your words
\end{itemize}
\vspace{0.3cm}

\begin{center}
\fcolorbox{orange!80}{orange!10}{%
\begin{minipage}{0.85\textwidth}
\centering
\textbf{Few-shot examples encode your judgment for edge cases}
\end{minipage}
}
\end{center}
\end{frame}

\section{More Techniques}

\begin{frame}{Changing Models and Role Specification}
\textbf{Not all models are equal for every task}:
\begin{itemize}
    \item \textbf{Larger models} (GPT-4, Claude Opus) -- better reasoning, more expensive
    \item \textbf{Smaller models} (Claude Haiku, GPT-4o-mini) -- faster, cheaper, good for simple tasks
    \item Start with a capable model, then try cheaper ones and see if quality holds
    \item Sometimes switching models fixes problems that prompt engineering can't
\end{itemize}
\vspace{0.3cm}

\textbf{Role/persona specification} -- set the model's ``mindset'':
\begin{itemize}
    \item ``You are a hiring manager at a tech company'' $\rightarrow$ evaluates like a practitioner
    \item ``You are a senior engineer reviewing technical qualifications'' $\rightarrow$ focuses on depth
    \item Activates relevant patterns from training data; one sentence is usually enough
\end{itemize}
\end{frame}

\begin{frame}{Chain-of-Thought and Output Validation}
\textbf{Chain-of-thought} -- ask the model to reason before answering:
\begin{itemize}
    \item ``First list the relevant qualifications, then evaluate each one''
    \item Forces organized reasoning before committing to an answer
    \item Tradeoff: more tokens (cost and latency) for better accuracy
\end{itemize}
\vspace{0.3cm}

\textbf{Output validation and guardrails} -- don't trust, verify:
\begin{itemize}
    \item \textbf{Schema validation}: Use structured output (e.g., Pydantic) to enforce field types and ranges
    \item \textbf{Range checks}: Score must be 0--100, not 150 or $-$5
    \item \textbf{Re-prompting}: If output fails validation, try again (with a limit)
    \item Combine AI flexibility with code reliability -- a safety net, not a replacement for good prompts
\end{itemize}
\end{frame}

\section{Beyond Prompt Engineering}

\begin{frame}{When Prompting Isn't Enough: Fine-Tuning}
\textbf{Fine-tuning} = training a model on your specific data
\vspace{0.3cm}

\textbf{Consider fine-tuning when}:
\begin{itemize}
    \item You have \emph{lots} of labeled examples
    \item Prompt engineering has plateaued and you need more consistency
    \item You need to reduce cost (fine-tuned small models can replace large ones)
    \item You have a very domain-specific task with specialized vocabulary
\end{itemize}
\vspace{0.3cm}

\textbf{Don't fine-tune when}:
\begin{itemize}
    \item You haven't tried prompt engineering thoroughly first
    \item Your labeled data is small or noisy
    \item Requirements change frequently (re-training is expensive)
\end{itemize}
\vspace{0.3cm}

\textbf{Rule of thumb}: Exhaust prompt engineering before fine-tuning.
\end{frame}

\section{Putting It All Together}

\begin{frame}{The Improvement Toolkit}
\begin{center}
\renewcommand{\arraystretch}{1.4}
\begin{tabular}{l l}
\textbf{Problem} & \textbf{Solution} \\
\hline
Scores bunch together & Domain context, anchor points \\
Model makes things up & Grounding with citations \\
Results vary across runs & Temperature, structured output \\
Misses important factors & Decomposition \\
Edge cases handled randomly & Few-shot examples \\
Wrong ``personality'' for task & Role/persona specification \\
Reasoning is shallow & Chain-of-thought \\
Bad outputs slip through & Output validation \\
Prompt engineering plateaus & Change model or fine-tune \\
\end{tabular}
\end{center}
\end{frame}

\begin{frame}{Order of Operations}
\textbf{When improving an AI system, try techniques in this order}:
\vspace{0.3cm}

\begin{enumerate}
    \item \textbf{Get specific} -- Add domain context, criteria, and anchor points
    \item \textbf{Add structure} -- Structured output, lower temperature
    \item \textbf{Ground it} -- Require citations and evidence
    \item \textbf{Show examples} -- Few-shot prompting for edge cases
    \item \textbf{Break it apart} -- Decompose into sub-tasks
    \item \textbf{Try a different model} -- Bigger, smaller, or specialized
    \item \textbf{Fine-tune} -- Only after exhausting the above
\end{enumerate}
\vspace{0.3cm}

\textbf{Each step is more effort than the last.} Start cheap and simple.
\end{frame}

\section{Your Turn}

\begin{frame}{Today's Exercise}
\textbf{Open the notebook}: \texttt{lecture\_3\_resume\_scorer\_improvement.ipynb}
\vspace{0.3cm}

\textbf{You have 6 resumes}:
\begin{itemize}
    \item 2 \textbf{gold} -- strong matches (should score high)
    \item 2 \textbf{silver} -- decent matches (should score moderately)
    \item 2 \textbf{wild} -- unknown quality (likely score lower)
\end{itemize}
\vspace{0.3cm}

\textbf{Your goal}:
\begin{itemize}
    \item Start with the baseline single-shot scorer
    \item Apply the techniques from today's lecture
    \item Improve separation: gold $>$ silver $>$ wild
    \item Submit your scores to the leaderboard
\end{itemize}
\end{frame}

\begin{frame}{What to Try}
\textbf{Pick techniques and experiment}:
\vspace{0.3cm}

\begin{itemize}
    \item Rewrite the prompt with specific scoring criteria
    \item Add a role/persona to the system message
    \item Include few-shot examples of good and bad resumes
    \item Decompose into extraction steps and combine in code
    \item Require evidence citations for the score justification
    \item Try different models and compare results
    \item Adjust temperature for more consistent outputs
\end{itemize}
\vspace{0.3cm}

\textbf{Measure your progress}: Do gold resumes reliably outscore silver? Do silver outscore wild?
\end{frame}

\begin{frame}{Let's Get Started!}
\begin{center}
\Large
Open the notebook: \\[0.3cm]
\texttt{lecture\_3\_resume\_scorer\_improvement.ipynb}
\\[1cm]

Apply the techniques \\[0.3cm]
Improve gold vs. silver vs. wild separation \\[0.3cm]
Submit to the leaderboard
\end{center}
\end{frame}

\end{document}
